% 第8章 8.2.2：E(x,y)=exp(xy) 在 [-1,1]^2 上的 r 项分离逼近，r=1,...,8（对数纵轴）。
% (a) Taylor 截断 E_r 的一致误差 e-sum_{k<r}1/k!（蓝）与上界 e/r!（蓝虚线）；
% (b) L2 误差：Taylor 截断（绿）与最佳秩 r 截断 sqrt(sum_{j>r} sigma_j^2)（蓝），sigma_j 为 Legendre 正交基系数矩阵的奇异值。
% 数值由同目录 generate.py 计算（48 点 Gauss-Legendre 求积，14 个基函数）。
\begin{tikzpicture}[font=\figurefont\footnotesize]
\pgfplotsset{every axis/.style={axis lines=left,
  x axis line style={draw=FigureStroke,line width=.6pt},y axis line style={draw=FigureStroke,line width=.6pt},
  tick style={draw=FigureStroke},tick label style={font=\figurefont\scriptsize,text=FigureInk},
  title style={font=\figurefont\footnotesize},label style={font=\figurefont\scriptsize},
  legend style={draw=none,fill=none,font=\figurefont\scriptsize,cells={anchor=west}}}}
\pgfplotsset{every axis/.append style={width=52mm,height=40mm,xmin=0.6,xmax=8.4,xtick={1,...,8},xlabel={项数 $r$},
  ymajorgrids,grid style={draw=FigureGrid,line width=.4pt}}}
\begin{semilogyaxis}[name=A,ymin=1e-5,ymax=10,title={(a) 一致误差},legend style={at={(0.03,0.03)},anchor=south west}]
\addplot[draw=FigureBlue,line width=1.1pt,mark=*,mark size=1.4pt,mark options={fill=FigureBlue,draw=FigureBlue}] coordinates {(1,1.718) (2,0.7183) (3,0.2183) (4,5.162e-2) (5,9.948e-3) (6,1.615e-3) (7,2.263e-4) (8,2.786e-5)};
\addplot[draw=FigureBlue,line width=1.1pt,dashed] coordinates {(1,2.718) (2,1.359) (3,0.4530) (4,0.1133) (5,2.265e-2) (6,3.775e-3) (7,5.393e-4) (8,6.742e-5)};
\legend{Taylor 截断,$e/r!$}
\end{semilogyaxis}
\begin{semilogyaxis}[at={(A.east)},anchor=west,xshift=14mm,ymin=5e-10,ymax=10,title={(b) $L^2$ 误差},legend style={at={(0.03,0.03)},anchor=south west}]
\addplot[draw=FigureGreen,line width=1.1pt,mark=square*,mark size=1.3pt,mark options={fill=FigureGreen,draw=FigureGreen}] coordinates {(1,0.7383) (2,0.2144) (3,4.999e-2) (4,9.593e-3) (5,1.556e-3) (6,2.182e-4) (7,2.690e-5) (8,2.958e-6)};
\addplot[draw=FigureBlue,line width=1.1pt,mark=*,mark size=1.4pt,mark options={fill=FigureBlue,draw=FigureBlue}] coordinates {(1,0.7138) (2,8.950e-2) (3,7.645e-3) (4,4.847e-4) (5,2.446e-5) (6,1.026e-6) (7,3.682e-8) (8,1.155e-9)};
\legend{Taylor 截断,最佳秩 $r$}
\end{semilogyaxis}
\end{tikzpicture}
